\DocumentMetadata{
  pdfversion=2.0,
  pdfstandard=ua-2,
  lang=en-US,
  tagging=on,
  tagging-setup={math/setup=mathml-SE,tabsorder=structure}
}
\documentclass[11pt,letterpaper]{article}
\usepackage[margin=0.68in,includefoot]{geometry}
\usepackage{newcomputermodern}
\usepackage{amsmath,amscd,mathtools}
\usepackage{tikz,adjustbox}
\providecommand{\square}{\mdlgwhtsquare}
\usepackage[hidelinks]{hyperref}
\hypersetup{
  pdftitle={Combinatorics PhD Exam, May 2011},
  pdfauthor={University of Florida Department of Mathematics},
  pdfsubject={Graduate mathematics examination},
  pdfdisplaydoctitle=true,
  bookmarksopen=true
}
\setcounter{secnumdepth}{0}
\setlength{\parindent}{0pt}
\setlength{\parskip}{0.28em}
\setlength{\emergencystretch}{3em}
\setlength{\leftmargini}{2.15em}
\setlength{\leftmarginii}{2.65em}
\setlength{\leftmarginiii}{3em}
\pagestyle{empty}
\raggedbottom
\allowdisplaybreaks
\NewTaggingSocketPlug{block/list/label}{examproblem}{%
  \tagstructbegin{tag=Lbl}\tagstructend%
  \tagstructbegin{tag=\LBody}%
  \tagstructbegin{tag=\ProblemHeadingTag,title-o={\ProblemHeadingText}}%
  \tagmcbegin{tag=\ProblemHeadingTag,actualtext={\ProblemHeadingText}}%
  #2%
  \tagmcend%
  \tagstructend%
}
\newenvironment{examproblems}[2]{%
  \def\ProblemHeadingTag{#2}%
  \AssignTaggingSocketPlug{block/list/label}{examproblem}%
  \begin{enumerate}%
  \setcounter{enumi}{#1}%
  \renewcommand{\labelenumi}{\textbf{\theenumi.}}%
  \setlength{\topsep}{0.35em}%
  \setlength{\partopsep}{0pt}%
  \setlength{\itemsep}{0.48em}%
  \setlength{\parsep}{0.18em}%
}{\end{enumerate}}
\newenvironment{examsubparts}{%
  \AssignTaggingSocketPlug{block/list/label}{default}%
  \begin{enumerate}%
  \setlength{\topsep}{0.18em}%
  \setlength{\partopsep}{0pt}%
  \setlength{\itemsep}{0.16em}%
  \setlength{\parsep}{0.1em}%
}{\end{enumerate}}
\newenvironment{examinstructions}{%
  \par\begingroup\itshape
}{%
  \par\endgroup\vspace{0.18em}
}
\makeatletter
\renewcommand\subsection{\@startsection{subsection}{2}{0pt}%
  {-1.25ex plus -.25ex minus -.1ex}%
  {0.55ex plus .1ex}%
  {\normalfont\large\bfseries}}
\renewcommand{\@maketitle}{%
  \newpage\null\vskip -1em%
  {\footnotesize\bfseries
    \href{https://www.ufl.edu/}{University of Florida}, \href{https://math.ufl.edu/}{Department of Mathematics}\par}%
  \vskip -0.5em%
  \tagstructbegin{tag=H1,title={Combinatorics PhD Exam, May 2011}}%
  \tagpdfparaOff%
  \tagmcbegin{tag=H1}%
    {\LARGE\bfseries\href{https://gma.math.ufl.edu/exams/combinatorics-phd/}{Combinatorics PhD Exam}, May 2011\par}%
  \tagmcend%
  \tagpdfparaOn%
  \tagstructend%
  \vskip 0.25em%
  \tagmcbegin{artifact}%
    \hrule height 1pt%
  \tagmcend%
  \vskip 1em%
}
\makeatother
\title{Combinatorics PhD Exam, May 2011}
\author{}
\date{}
\begin{document}
\maketitle
\thispagestyle{empty}
\begin{examproblems}{0}{H2}
\def\ProblemHeadingText{Problem 1.}
\item
Let \(f(n)\) be the number of permutations of length~\(n\) that have no cycle that is shorter than three.
\begin{examsubparts}
\item[\textnormal{(a)}]
Find the exponential generating function of the sequence \(f(n)\). You can assume that \(f(0)=1\).
\item[\textnormal{(b)}]
Find \(\lim_{n\to\infty} f(n)/n!\).
\end{examsubparts}
\def\ProblemHeadingText{Problem 2.}
\item
We divide a group of~\(n\) people into subgroups~\(A\), \(B\), and~\(C\), and ask each subgroup to form a line. We also require that~\(A\) have an odd number of people, and that~\(B\) have an even number of people. How many ways are there to do this?
\def\ProblemHeadingText{Problem 3.}
\item
Let~\(A\) be the graph obtained from \(K_n\) by deleting an edge. Find a formula for the number of spanning trees of~\(A\).
\def\ProblemHeadingText{Problem 4.}
\item
Let \(X_n(p)\) be the number of cycles of the~\(n\)-permutation~\(p\). Compute \(\operatorname{VAR}(X_n)\).
\def\ProblemHeadingText{Problem 5.}
\item
Show that a planar graph for which every face has an even number of edges must be bipartite.
\def\ProblemHeadingText{Problem 6.}
\item
Recall that a tournament is a complete directed graph. Prove that there exists a tournament on eight vertices that contains at least 316 Hamiltonian paths.
\def\ProblemHeadingText{Problem 7.}
\item
For which values of~\(m\) and~\(n\) is the complete bipartite graph \(K_{m,n}\) planar, for which values of~\(m\) and~\(n\) is it Hamiltonian, and for which values of~\(m\) and~\(n\) is it Eulerian?
\def\ProblemHeadingText{Problem 8.}
\item
Prove in a finite poset~\(P\), the number of elements in the longest chain is equal to the smallest number~\(k\) so that~\(P\) can be decomposed into the union of~\(k\) antichains.
\end{examproblems}
\end{document}
